\chapter{Der Märtyrer Euplius}
\quellseite{8}

\begin{quote}
„Will mir jemand nachfolgen, der verleugne sich selbst, und nehme sein Kreuz auf sich und folge mir.“
\end{quote}
\quelle{Matth. 16, 24.}

Am 12. August des Jahres 304 hatte sich zu Catana in Sicilien der römische Statthalter Calvisian zu Gericht gesetzt. Hinter dem Richterstuhl war nach römischer Sitte ein Vorhang heruntergelassen, damit sich die Richter ungestört zurückziehen, berathschlagen und das Urtheil sprechen könnten. Calvisian befand sich noch hinter dem Vorhang, als Cuplius in das Gerichtszimmer geführt wurde. Der Statthalter bestieg den Richterstuhl. Cuplius trat vor ihn, das Evangelienbuch in der Hand. Zornig rief ihm Maximus, ein Freund des Statthalters, zu: „Wie kannst du es wagen, hier mit einem Buch zu erscheinen, das nach den Geboten der Kaiser ausgeliefert und verbrannt werden soll?“ Und Calvisian fragte: „Woher hast du diese
\quellseite{9}
Schriften?“ „Bringst du sie mit aus deinem Hause?“ Cuplius entgegnete: „Ach nein, ich habe kein Haus; ich hatte das Buch bei mir, als ich verhaftet wurde.“

Der Richter gebot ihm nun, etwas aus dem Evangelienbuch zu lesen. Cuplius schlug Matth. 5, 10. auf und las: „Selig sind, die um Gerechtigkeit willen verfolgt werden, denn das Himmelreich ist ihr.“ Dann schlug er einige Blätter um und las weiter Matth. 16, 24.: „Will mir jemand nachfolgen, der verleugne sich selbst und nehme sein Kreuz auf sich und folge mir.“ „Dies ist“, fügte er hinzu, „das Gebot, das mein Gott mir gegeben hat.“ „Durch wen?“ fragte der Römer. „Durch Jesum Christum, den Sohn des lebendigen Gottes“, antwortete der Christ. Hierauf rief Salvian: „Da Cuplius selbst gesteht, daß er ein Christ ist, so soll er dem Schergen übergeben und auf die Folter gespannt werden.“ Sein Befehl wurde sofort auf das grausamste vollzogen.

Nach einiger Zeit hieß der Richter mit der Peinigung inne halten und dem Märtyrer ein paar Stunden Ruhe gewähren. Aber nicht aus Mitleid. Die Schmerzen sollten in dieser Frist nur besser nachwirken und den Märtyrer mürbe machen. Nach Verlauf dieser Frist wurde er zum zweiten Male auf die Marterbank gespannt und während der neuen heftigen Qualen wieder gefragt, ob er jetzt noch immer bei seiner Meinung bleiben wolle. Cuplius antwortete: „Ich habe es schon bekannt und bekenne es aufs neue, daß ich ein Christ bin und die heiligen Schriften lese. Wollte ich sie ausliefern oder verleugnen, so würde ich meinen Gott beleidigen; darum will ich lieber sterben, als einen solchen Frevel begehen; denn auf meinen Tod wird ein glückseliges Leben folgen.“ Da befahl Salvian, die Qualen zu verdoppeln. Cuplius aber betete laut: „Ich danke dir, o Herr, daß du mich würdigst, um deines Namens willen zu leiden. Darum errettete mich! Das ist es, worum ich dich anflehe.“ Salvian rief ihm zu: „Entsage deiner Thorheit! Bete unsere Götter an, so will ich dich in Freiheit setzen.“ Aber der Christ entgegnete: „Ich bete Jesum Christum an und verabscheue die Teufel. Ich bin ein Christ, so verdamme mich nun, wenn du willst, zu neuen Qualen; ich will meinem Erlöser zu Lieb diese Leiden erdulden.“

Wieder ließ der erzürnte Statthalter mit der Peinigung inne halten. Cuplius wurde auf seinen Befehl aus der Marterbank gespannt, und er schrie ihm nochmals zu: „Bete unsere Götter an, du Bösewicht! Bete den Mars, den Apollo, den Æsculap an.“ Cuplius aber rief laut: „Ich bete den Vater, den Sohn und den Heiligen Geist an; ich bete den dreieinigen Gott an, denn er ist der allein wahre Gott, und es gibt keinen andern außer ihm.“ Da loderte Salvians Zorn aufs neue auf, und er befahl, die Folter zum dritten Male anzuwenden, und zwar mit aller nur erdenklichen Grausamkeit. Cuplius aber stärkte sich mit neuem Gebet. Unter den gräßlichsten Qualen rief er zu wiederholten Malen: „Ich danke dir, mein Gott, daß du mich solcher Leiden um deines heiligen Evangeliums willen würdigest; mein Jesu, stehe mir bei; um deines Namens willen leide ich diese Marter!“ Und als unter den furchtbaren Schmerzen seine Kräfte schon so sehr geschwunden waren, daß er nicht mehr laut sprechen konnte, bewegten sich immer noch leise seine Lippen.

Nun endlich trat Salvian hinter den Vorhang, um das Urtheil zu fällen. Nach einer Weile schritt er wieder hervor und verlas folgenden Richterspruch: „Wir befehlen, daß Cuplius, den wir überführt haben, daß er ein Christ ist, zur Strafe für seine Widerspenstigkeit gegen die Gesetze der Kaiser und gegen die Götter, enthauptet werde.“

Hierauf wurde dem Märtyrer das Evangelienbuch um den Hals gebunden. Ein öffentlicher Ausrufer ging vor ihm her und rief mit lauter Stimme: „Sehet hier Cuplius, den Christen und den Feind der Kaiser und der Götter!“ Auf seinem Gang zum Richtplatz rief

der Verurtheilte den Hörern an: „Herr Jesu Christe, ich danke dir! Stärke nun auch, was du in mir gewirkt hast!“ Auf der Richtstätte betete er noch lange auf seinen Knieen; dann bot er dem Henker sein Haupt willig zum Todesstreich dar und ging ein in die Wohnungen der ewigen Seligkeit.
